\paragraph{\trYonedaEmbedding。}\label{sub:yoneda-embedding}
对 $\scr{A}$ 的每个\trObject{} $Y$，都有一个相伴的\trNonunital{} \trAinfty-模
$\scr{Y}$，由 $\scr{Y}(X)=\hom_{\scr{A}}(X,Y)$ 与 $\mu^d_{\scr{Y}}=\mu^d_{\scr{A}}$
给出定义。这是一个\trCanonical{}的\trNonunital{} \trAFunctor
$\scr{J}=\scr{J}_{\scr{A}}\colon\scr{A}\to\scr{Q}
=\mathit{nu}\text{-}\mathit{mod}(\scr{A})$ 的一部分，该\trFunctor{}称为%
\textbf{\trYonedaEmbedding}（Yoneda embedding）。设
$c\in\hom_{\scr{A}}(Y_0,Y_1)$，则
$\scr{J}^1(c)\in\hom_{\scr{A}}(\scr{Y}_0,\scr{Y}_1)$\footnote{译注：原书此处
印作 $\hom_{\scr{A}}(\scr{Y}_0,\scr{Y}_1)$；疑应为
$\hom_{\scr{Q}}(\scr{Y}_0,\scr{Y}_1)$（因 $\scr{Y}_0$、$\scr{Y}_1$ 是
$\scr{Q}=\mathit{nu}\text{-}\mathit{mod}(\scr{A})$ 的\trObject{}，而非 $\scr{A}$ 的%
\trObject）。此处照原书排印，存疑备案。——译者注}就是预\trModule\trHomomorphism
\begin{equation*}
\begin{aligned}
&\scr{Y}_0(X_{d-1})\otimes\hom_{\scr{A}}(X_{d-2},X_{d-1})
\otimes\cdots\otimes\hom_{\scr{A}}(X_0,X_1)
\longrightarrow\scr{Y}_1(X_0),\\
&(b,a_{d-1},\ldots,a_1)\longmapsto\mu^{d+1}_{\scr{A}}(c,b,a_{d-1},\ldots,a_1);
\end{aligned}
\end{equation*}
高阶项 $\scr{J}^d$（$d\ge 2$）则由这个公式的显然类似物来定义。

\trYonedaEmbedding{}与\trPullback\trCompatible，其确切含义如下。设
$\scr{F}\colon\scr{A}\to\scr{B}$
为一个\trNonunital{} \trAFunctor。那么存在一个\trCanonical{}\trNaturalTransformation
\begin{equation}\label{eq:yoneda-pullback-transformation}
T\colon\scr{J}_{\scr{A}}\longrightarrow\scr{F}^*\circ\scr{J}_{\scr{B}}\circ\scr{F}.
\end{equation}
一阶项 $T^1$ 由对每个 $Y$ 的一个\trModule\trHomomorphism{}组成，而这个%
\trModule\trHomomorphism{}又是一系列映射
\begin{equation}\label{eq:yoneda-first-order-maps}
\begin{split}
&(\scr{J}_{\scr{A}}(Y))(X_{d-1})\otimes\hom_{\scr{A}}(X_{d-2},X_{d-1})\\
&\quad\otimes\cdots\otimes\hom_{\scr{A}}(X_0,X_1)\\
&\qquad\longrightarrow(\scr{F}^*\circ\scr{J}_{\scr{B}}\circ\scr{F}(Y))(X_0).
\end{split}
\end{equation}
\begingroup\emergencystretch=8em\hbadness=10000
按定义\footnote{译注：原书\eqref{eq:yoneda-first-order-maps} 末行印作
$\scr{F}(Y))(X_0)$（$(Y)$ 在外层括号之内）；疑应为 $\scr{F})(Y)(X_0)$
（与本句按定义式及映射的类型相合；同页下文按定义句即如此排印），
显见排印笔误；此处照原书排印，存疑备案。——译者注}，
$(\scr{J}_{\scr{A}}(Y))(X_{d-1})=\hom_{\scr{A}}(X_{d-1},Y)$，而
$(\scr{F}^*\circ\scr{J}_{\scr{B}}\circ\scr{F})(Y)(X_0)
=\hom_{\scr{B}}(\scr{F}(X_0),\scr{F}(Y))$\footnote{译注：原书此处下标印作斜体
$B$（即 $\mathcal{J}_B$）；疑应为花体 $\scr{B}$（与\eqref{eq:yoneda-pullback-transformation}
中 $\scr{J}_{\scr{B}}$ 一致），显见排印笔误；此处径改。——译者注}；于是，
\eqref{eq:yoneda-first-order-maps} 就是映射
$(b,a_{d-1},\ldots,a_1)\longmapsto(-1)^{\clubsuit_{d-1}+|b|}\scr{F}^d(b,a_{d-1},\ldots,a_1)$。
\par\endgroup
类似的表达式定义 $T^r$（$r\ge 2$）。

换个方向，可以把 $\scr{J}^1$ 推广成一个\trChainMap
\begin{equation}\label{eq:yoneda-evaluation-chain-map}
\begin{split}
\lambda&=\lambda_{\scr{M}}\colon\scr{M}(Y)
\longrightarrow\hom_{\scr{Q}}(\scr{Y},\scr{M}),\\
\lambda(c)^d(b,a_{d-1},\ldots,a_1)&=\mu^{d+1}_{\scr{M}}(c,b,a_{d-1},\ldots,a_1)
\end{split}
\end{equation}
对任意\trNonunital{} \trAinfty-模 $\scr{M}$。在\trCohomology{}层面上，关于在任一侧作
乘法，它有如下易于验证的性质：
\begin{equation}\label{eq:lambda-multiplication-property}
[\lambda_{\scr{M}}(\mu^2_{\scr{M}}(c,b))]
=[\mu^2_{\scr{Q}}(\lambda_{\scr{M}}(c),\scr{J}^1(b))]
\in\Hom_{H(\scr{Q})}(\scr{Y}_0,\scr{M})
\end{equation}
其中 $[c]\in H(\scr{M}(Y_1))$、$[b]\in\Hom_{H(\scr{A})}(Y_0,Y_1)$；以及
\begin{equation}\label{eq:lambda-second-property}
[\mu^2_{\scr{Q}}(t,\lambda_{\scr{M}_0}(c))]
=[\lambda_{\scr{M}_1}(t^1(c))]
\in\Hom_{H(\scr{Q})}(\scr{Y},\scr{M}_1)
\end{equation}
其中 $[t]\in\Hom_{H(\scr{Q})}(\scr{M}_0,\scr{M}_1)$、$[c]\in H(\scr{M}_0(Y))$。注意，
如果我们只限于那些来自\trYonedaEmbedding{}的\trModule，并相应地把 $t$ 取成 $\scr{A}$
中某个\trMorphism{}的像，那么这两个方程都归结为 $\scr{J}$ 的\trAFunctor{}方程的特例。

\begin{remark}\label{thm:linearization-functor-picture}
映射 $\lambda$ 可以解释为一幅更大图景的一部分，而这幅图景同时也解释了上面列出的那些
性质。定义\textbf{\trLinearization\trFunctor}（linearization functor）为下述%
\trNonunital{} dg \trFunctor：它把 $\scr{Q}$ 的\trYonedaEmbedding{}与用 $\scr{A}$ 的
那一个所作的\trPullback{}结合起来，
\begin{equation*}
\mathscr{L}\colon\scr{Q}\xrightarrow{\;\scr{J}_{\scr{Q}}\;}
\mathit{nu}\text{-}\mathit{mod}(\scr{Q})\xrightarrow{\;\scr{J}^*_{\scr{A}}\;}\scr{Q}
\end{equation*}
（这个名称的由来是：每个\trModule\trHomomorphism{}在 $\mathscr{L}$ 下的像都是一个%
\trHomomorphism，其唯一非零的项就是一阶的那一项，即\trLinear{}的那一项）。这个%
\trFunctor{}附带一个\trCanonical{}\trNaturalTransformation{}%
$L\colon\mathrm{Id}_{\scr{Q}}\to\mathscr{L}$。在一阶上，它由一个\trModule\trHomomorphism%
{}给出，即 $l=L^1_{\scr{M}}\colon\scr{M}\to\scr{M}_{\mathrm{lin}}=\mathscr{L}(\scr{M})$
（对每个 \trAinfty-模 $\scr{M}$），而其起始项由\trLinear{}映射
$l^1\colon\scr{M}(Y)\to\scr{M}_{\mathrm{lin}}(Y)$（对一切 $Y$）组成。按定义，
$\scr{M}_{\mathrm{lin}}(Y)=\hom_{\scr{Q}}(\scr{Y},\scr{M})$，其中 $\scr{Y}$ 是 $Y$ 在%
\trYonedaEmbedding{}下的像，而这个\trLinear{}映射就是我们的 $\lambda$（至多差一个正负号
$(-1)^{|c|}$）。$l$ 是 \trAinfty-模之间的\trHomomorphism{}这一事实蕴涵
\eqref{eq:lambda-multiplication-property}；而 $L$ 是\trNaturalTransformation{}这一事实
蕴涵\eqref{eq:lambda-second-property}。
\end{remark}
